\documentclass[a4paper,11pt]{article}

% Packages
\usepackage[utf8]{inputenc}
\usepackage[T1]{fontenc}
\usepackage{lmodern}
\usepackage{microtype}
\usepackage[french]{babel}
\usepackage{geometry}
\usepackage{xcolor}
\usepackage{enumitem}
\usepackage{titlesec}
\usepackage{tcolorbox}
\usepackage{graphicx}
\usepackage{fancyhdr}
\usepackage{tabularx}
\usepackage{booktabs}
\usepackage{multicol}
\usepackage{tikz}
\usepackage{amsmath}
\usepackage{amssymb}
\usepackage{hyperref}
\usepackage{listings}
\usepackage{fontawesome5}

% ============================================================
% Couleurs violet/lavande
% ============================================================
\definecolor{mainpurple}{RGB}{106, 90, 205}
\definecolor{lightlavender}{RGB}{230, 230, 250}
\definecolor{deeppurple}{RGB}{75, 0, 130}
\definecolor{softpurple}{RGB}{147, 112, 219}
\definecolor{palelavender}{RGB}{245, 240, 255}

% Couleurs code (thème sombre)
\definecolor{codebg}{RGB}{30, 30, 30}
\definecolor{codegreen}{RGB}{0, 200, 100}
\definecolor{codegray}{RGB}{160, 160, 160}
\definecolor{codestring}{RGB}{206, 145, 120}
\definecolor{codekeyword}{RGB}{86, 156, 214}
\definecolor{codecomment}{RGB}{106, 153, 85}

% ============================================================
% Géométrie
% ============================================================
\geometry{
    top=2.5cm,
    bottom=2.5cm,
    left=2.5cm,
    right=2.5cm,
    headheight=22pt
}

% ============================================================
% En-têtes / pieds de page
% ============================================================
\pagestyle{fancy}
\fancyhf{}
\fancyhead[L]{\color{mainpurple}\small\leftmark}
\fancyhead[R]{\color{softpurple}\small PROJET FTP - SR}
\fancyfoot[C]{\color{mainpurple}\thepage}
\renewcommand{\headrulewidth}{0.5pt}
\renewcommand{\footrulewidth}{0.5pt}
\renewcommand{\headrule}{\hbox to\headwidth{\color{mainpurple}\leaders\hrule height \headrulewidth\hfill}}
\renewcommand{\footrule}{\hbox to\headwidth{\color{mainpurple}\leaders\hrule height \footrulewidth\hfill}}

% ============================================================
% Titres
% ============================================================
\titleformat{\section}
  {\LARGE\bfseries\color{deeppurple}}
  {\thesection}{1em}{}
  [\titlerule]

\titleformat{\subsection}
  {\Large\bfseries\color{mainpurple}}
  {\thesubsection}{1em}{}

\titleformat{\subsubsection}
  {\large\bfseries\color{softpurple}}
  {\thesubsubsection}{1em}{}

% ============================================================
% Boîtes tcolorbox (uniquement celles du document)
% ============================================================
\newtcolorbox{infobox}{
    colback=lightlavender,
    colframe=mainpurple,
    boxrule=1pt,
    arc=3mm,
    left=10pt, right=10pt, top=10pt, bottom=10pt
}

\newtcolorbox{tipbox}{
    colback=palelavender,
    colframe=softpurple,
    boxrule=1pt,
    arc=3mm,
    left=10pt, right=10pt, top=10pt, bottom=10pt,
    title={\textbf{$\triangleright$ Idée}}
}

\newtcolorbox{warningbox}{
    colback=lightlavender!50,
    colframe=deeppurple,
    boxrule=1.5pt,
    arc=3mm,
    left=10pt, right=10pt, top=10pt, bottom=10pt,
    title={\textbf{$\blacktriangledown$ Attention}}
}

\newtcolorbox{examplebox}{
    colback=palelavender,
    colframe=mainpurple,
    boxrule=1pt,
    arc=2mm,
    left=8pt, right=8pt, top=8pt, bottom=8pt,
    title={\textbf{$\star$ Exemple concret}}
}

\newtcolorbox{sciencebox}{
    colback=lightlavender!70,
    colframe=deeppurple,
    boxrule=1pt,
    arc=3mm,
    left=10pt, right=10pt, top=10pt, bottom=10pt,
    title={\textbf{$\boldsymbol{\Sigma}$ Éclairage technique}}
}

% ============================================================
% Listings (blocs de code)
% ============================================================
\lstdefinestyle{codestyle}{
    backgroundcolor=\color{white},
    basicstyle=\ttfamily\small\color{black},
    keywordstyle=\color{mainpurple}\bfseries,
    stringstyle=\color{red!85!orange},
    commentstyle=\color{deeppurple}\itshape,
    numberstyle=\tiny\color{red!85!orange},
    breaklines=true,
    breakatwhitespace=false,
    keepspaces=true,
    showspaces=false,
    showstringspaces=false,
    frame=single,
    framesep=4pt,
    rulecolor=\color{mainpurple},
    xleftmargin=8pt,
    xrightmargin=4pt,
    tabsize=4,
    literate=
        {←}{{\textcolor{codegreen}{<-}}}2
        {→}{{\textcolor{codegreen}{->}}}2
        {✅}{{\textcolor{codegreen}{[OK]}}}3
        {❌}{{\textcolor{red}{[NON]}}}4
}

\lstdefinestyle{bash}{
    style=codestyle,
    language=bash,
    morekeywords={adb,nmap,grep,sed,zip,sudo,apt}
}

\lstdefinestyle{cstyle}{
    style=codestyle,
    language=C
}

\lstset{style=codestyle}

% ============================================================
% Liens
% ============================================================
\hypersetup{
    colorlinks=true,
    linkcolor=deeppurple,
    citecolor=mainpurple,
    urlcolor=softpurple,
    pdfborder={0 0 0},
    pdftitle={TP FTP -- Documentation},
    pdfauthor={ANGERETTI Quentin et CASSOU Baptiste}
}

% ============================================================
% Informations de garde
% ============================================================
\newcommand{\titreprojet}{TP FTP}
\newcommand{\soustitreprojet}{Systèmes et Réseaux}
\newcommand{\nometudiantA}{ANGERETTI Quentin}
\newcommand{\nometudiantB}{CASSOU Baptiste}
\newcommand{\anneeuniversitaire}{2025--2026}

\begin{document}

% ============================================================
% PAGE DE TITRE
% ============================================================
\begin{titlepage}
    \centering
    \vspace*{1.2cm}

    {\Huge\bfseries\color{deeppurple} \titreprojet\par}
    \vspace{0.3cm}
    {\LARGE\bfseries\color{mainpurple} \soustitreprojet\par}

    \vspace{1.2cm}

    \begin{tcolorbox}[
        colback=lightlavender,
        colframe=mainpurple,
        width=0.9\textwidth,
        arc=5mm,
        boxrule=2pt
    ]
    \centering
    {\Large\color{deeppurple} TP : Transfert de fichiers (FTP)\par}
    \vspace{0.35cm}
    {\large Service de transfert de fichiers inspiré du protocole FTP\par}
    \vspace{0.2cm}
    {\normalsize Client, serveur concurrent, reprise de transfert, architecture maître/esclaves\par}
    \end{tcolorbox}

    \vspace{1.1cm}

    \begin{tikzpicture}
        \fill[mainpurple!30] (0,0) circle (2.8cm);
        \fill[softpurple!40] (0,0) circle (1.9cm);
        \fill[deeppurple!20] (0,0) circle (1.05cm);
        \node[align=center] at (0,0) {\color{deeppurple}\bfseries FTP};
    \end{tikzpicture}

    \vfill

    \begin{tcolorbox}[
        colback=palelavender,
        colframe=softpurple,
        width=0.75\textwidth,
        arc=4mm,
        boxrule=1.2pt
    ]
    \centering
    {\large\textbf{Réalisé par}\par}
    \vspace{0.25cm}
    {\large \nometudiantA\par}
    {\large \nometudiantB\par}
    \vspace{0.35cm}
    {\normalsize Année universitaire : \anneeuniversitaire\par}
    {\normalsize \today\par}
    \end{tcolorbox}
\end{titlepage}

\tableofcontents
\newpage

% ============================================================
\section{Analyse du sujet}

La vraie difficulté est la \textbf{conception du protocole} :
quoi envoyer, dans quel ordre, et comment gérer les erreurs.
\begin{infobox}
Le sujet demande de construire un service de transfert de fichiers inspiré de FTP, en plusieurs étapes, du cas simple au cas distribué.
Ce TP mélange :
\begin{itemize}[leftmargin=1.2cm]
    \item communication réseau (sockets TCP),
    \item protocole applicatif (messages client/serveur),
    \item gestion de fichiers (texte et binaire),
    \item concurrence et robustesse.
\end{itemize}
\end{infobox}


\subsection{FTP : rappel simple}
Le FTP est un protocol de transfert de fichiers :
\begin{sciencebox}
\textbf{FTP (File Transfer Protocol)} : un client envoie une commande (\texttt{get}, \texttt{put}, \texttt{ls}...), le serveur répond avec un code de statut et éventuellement des données.
TCP transporte des octets de façon fiable, mais ne donne pas le sens des messages : ce sens est défini par \textbf{votre protocole applicatif}.
\end{sciencebox}


\subsection{But global du projet}
Dans la version minimale, le système comporte :
\begin{itemize}[leftmargin=1.2cm]
    \item un client qui envoie une requête ;
    \item un serveur qui traite la requête et répond.
\end{itemize}

\begin{examplebox}
\textbf{Exemple.} Le client demande \texttt{\textcolor{softpurple}{get} rapport.pdf}. Le serveur renvoie le fichier s'il existe, sinon un code d'erreur.
\end{examplebox}

Le scénario de base :
\begin{enumerate}[leftmargin=1.2cm]
    \item connexion TCP client \textrightarrow{} serveur ;
    \item envoi d'une requête structurée ;
    \item traitement serveur ;
    \item réponse (succès/erreur + données).
\end{enumerate}

\subsection{Notion centrale : protocole applicatif}
\begin{tipbox}
Définir dès le début \texttt{typereq\_t}, \texttt{request\_t} et \texttt{response\_t}. Le protocole doit rester extensible (\texttt{bye}, \texttt{ls}, \texttt{put}, \texttt{rm}).
\end{tipbox}

\section{Plan d'action}
\subsection{Serveur FTP de base}
\begin{infobox}
Objectif : un serveur concurrent (pool de processus), port 2121, qui traite \texttt{GET} et ferme proprement sur \texttt{SIGINT}.
\end{infobox}

\textbf{Points clés :}
\begin{itemize}[leftmargin=1.2cm]
    \item un fichier par connexion (dans cette étape),
    \item codes de retour explicites côté client/serveur,
    \item support des fichiers binaires (raisonner en octets et tailles, pas en chaînes C),
    \item arrêt propre du serveur et nettoyage des processus fils.
\end{itemize}

\subsection{Amélioration et robustesse}
\begin{infobox}
Objectif : ne plus charger un gros fichier d'un coup, mais l'envoyer \textbf{par blocs}, puis permettre plusieurs commandes sur une même connexion.
\end{infobox}

\begin{itemize}[leftmargin=1.2cm]
    \item définir la taille de bloc et le mécanisme de fin de transfert ;
    \item conserver la connexion ouverte jusqu'à \texttt{bye} ;
    \item ajouter la reprise de transfert après interruption (offset).
\end{itemize}

\begin{warningbox}
Sans protocole clair sur les tailles et les offsets, la reprise peut corrompre le fichier (doublons, trous, troncature).
\end{warningbox}

\subsection{Idée-clé de maître / esclaves}

\textbf{Objectif} : répartir la charge. 

Un serveur maître reçoit le client, choisit un esclave (Round-Robin), puis le client parle directement à cet esclave.


\begin{tipbox}
Le sujet suppose des collections de fichiers identiques sur les esclaves, ce qui simplifie la répartition.
\end{tipbox}


\begin{examplebox}
Séquence type : 
\begin{center}
    
Client \textrightarrow{} Maître (demande) \textrightarrow{} Client (IP/port esclave) \textrightarrow{} Esclave (session FTP).

\end{center}
\end{examplebox}


\subsection{Commandes avancées et sécurité}
\begin{infobox}
Ajouter \texttt{ls}, \texttt{put} et \texttt{rm}, puis protéger les commandes sensibles avec authentification.
\end{infobox}

\begin{sciencebox}
La commande \texttt{ls} illustre l'interface système (\texttt{fork}/\texttt{exec}/\texttt{dup2} ou \texttt{popen}).

\texttt{put}/\texttt{rm} introduisent la mise à jour des fichiers et la question de cohérence entre esclaves.
\end{sciencebox}



\newpage
\section{Étape I : serveur FTP de base}

\subsection{Réponses directes aux questions Q1 à Q4}
\subsubsection*{Q2 -- Structure \texttt{request\_t}}
\begin{infobox}
La structure \texttt{request\_t} est définie dans \texttt{include/ftp\_shared.h}.  
Elle contient le type de requête (\texttt{type}) et le nom de fichier (\texttt{filename}), ainsi que des champs prévus pour les étapes suivantes (\texttt{version}, \texttt{offset}, \texttt{block\_size}, \texttt{login}, \texttt{password}).
\end{infobox}

\subsubsection*{Q3 -- Squelette client/serveur FTP}
\begin{infobox}
Le squelette 'echo' a été adapté et renommé en \texttt{clientFTP} / \texttt{serverFTP}.  
Le port est fixé à \texttt{2121} via \texttt{FTP\_PORT} (pas d'argument de port).  

La concurrence est assurée par un pool de processus de taille \texttt{NB\_PROC}.
\end{infobox}

\begin{examplebox}
\textbf{Validation de lancement :}
\begin{lstlisting}[style=bash]
make all
make run-server
make run-client HOST=127.0.0.1
\end{lstlisting}
\end{examplebox}

\subsubsection*{Q4 -- Terminaison propre avec \texttt{SIGINT}}
\begin{infobox}
Le parent installe un handler \texttt{SIGINT} qui :
\begin{itemize}[leftmargin=1.2cm]
    \item ferme la socket d'écoute,
    \item retransmet \texttt{SIGINT} à tous les fils du pool,
    \item attend la fin des fils avec \texttt{waitpid}.
\end{itemize}
Chaque fils installe aussi un handler \texttt{SIGINT} pour quitter proprement sa boucle d'acceptation.
\end{infobox}

\subsection{Implémentation des types partagés}
\begin{infobox}
Les types partagés sont centralisés dans \texttt{include/ftp\_shared.h}. Ils servent de contrat unique entre client, serveur, maître et esclaves.
\end{infobox}

\textbf{Extrait de la structure de requête :}
\begin{lstlisting}[style=cstyle]
typedef struct {
    uint32_t version;
    uint32_t type;      /* typereq_t */
    uint64_t offset;    /* reprise de transfert */
    uint32_t block_size;
    char filename[FTP_MAX_FILENAME];
    char login[FTP_MAX_LOGIN];
    char password[FTP_MAX_PASSWORD];
} request_t;
\end{lstlisting}

\begin{tipbox}
Cette base couvre déjà les étapes suivantes :
\texttt{GET} (étape I), transfert par blocs et reprise (\texttt{offset}, étape II), puis \texttt{ls}/\texttt{put}/\texttt{rm}/\texttt{auth} (étape IV), sans casser le format des messages.
\end{tipbox}

\subsection{Décisions de conception}
\begin{itemize}[leftmargin=1.2cm]
    \item \textbf{Versionnage} : le champ \texttt{version} permet de faire évoluer le protocole.
    \item \textbf{Réponse unifiée} : \texttt{response\_t} expose un \texttt{status} explicite pour la gestion d'erreurs.
    \item \textbf{Reprise native} : le champ \texttt{offset} est prévu dès la base.
    \item \textbf{Extensibilité} : les commandes futures sont déjà codées dans \texttt{typereq\_t}.
\end{itemize}

\subsection{Fonctionnement actuel de \texttt{serverFTP}}
\begin{infobox}
Le programme \texttt{serverFTP} ouvre une socket d'écoute sur \texttt{2121}, puis crée au démarrage un pool de \texttt{NB\_PROC} fils.
Chaque fils accepte des connexions et traite les clients ; le parent supervise le pool.
\end{infobox}

\begin{sciencebox}
Chaîne d'exécution côté serveur :
\begin{center}
\texttt{Open\_listenfd(2121) \textrightarrow{} Fork x NB\_PROC \textrightarrow{} Accept (dans les fils) \textrightarrow{} R/W socket}
\end{center}
\end{sciencebox}

\subsection{Fonctionnement actuel de \texttt{clientFTP}}
\begin{infobox}
Le programme \texttt{clientFTP} prend l'hôte en argument, se connecte au serveur sur le port \texttt{2121}, lit l'entrée standard et envoie les lignes au serveur.
Chaque ligne reçue en retour est affichée sur la sortie standard.
\end{infobox}

\textbf{Lancement local :}
\begin{lstlisting}[style=bash]
make all
make run-server
make run-client HOST=127.0.0.1
\end{lstlisting}

\subsection{Déroulé d'une session (version courante)}
\begin{enumerate}[leftmargin=1.2cm]
    \item le serveur démarre et écoute sur \texttt{2121},
    \item le client se connecte à \texttt{<host>:2121},
    \item l'utilisateur saisit une ligne dans le client,
    \item le serveur renvoie actuellement la même ligne (comportement de base de validation réseau),
    \item la session se termine à la fermeture du client.
\end{enumerate}


\end{document}
